\documentclass[12pt,a4paper]{article}

\usepackage[utf8]{inputenc}
\usepackage[T2A]{fontenc}
\usepackage[serbianc]{babel} 
\usepackage{amsmath, amsthm, amssymb} 
\usepackage{graphicx}
\usepackage{wrapfig} 
\usepackage{xcolor} 
\usepackage{geometry}
\geometry{a4paper, margin=2.5cm}
\usepackage{hyperref}

\hypersetup{
    colorlinks=true,
    linkcolor=black,
    filecolor=black,      
    urlcolor=blue,
    citecolor=black,
    hidelinks
}

\newtheorem{teorema}{Теорема}[section]
\newtheorem{definicija}{Дефиниција}[section]
\newtheorem{lema}{Лема}[section]

\title{\textbf{Основе, историјски развој и примена виртуелне реалности (VR)}}
\author{Александар Милошевић}
\date{\today}

\begin{document}

\maketitle

\tableofcontents
\vspace{1cm}
\hrule
\vspace{1cm}

\section{Увод и историјски развој виртуелне реалности}

Виртуелна реалност (VR) представља један од најфасцинантнијих технолошких прагова савременог доба. \textbf{Основна идеја} овог концепта јесте потпуно уклањање физичке стварности из видика корисника и његова замјена компјутерски генерисаним, интерактивним тродимензионалним окружењем. 

\begin{wrapfigure}{l}{0.38\textwidth}
  \centering
  \includegraphics[width=0.36\textwidth]{vr_slika.png}
  \caption{Приказ савременог VR система}
\end{wrapfigure}

Историја виртуелне реалности сеже деценијама уназад, много прије него што су моћни графички процесори постали доступни масама. \textit{Пионирски кораци} започели су још педесетих и шездесетих година 20. вијека. Мортон Хеилиг је 1962. године конструисао уређај под називом \textbf{Sensorama}, који је представљао механички кабинет за мултимедијално искуство – укључујући стереоскопски екран, вентилаторе за дочаравање вјетра, мирисе и стерео звук. Међутим, прави темељ за хардверски засноване HMD (Head-Mounted Display) системе поставио је Иван Садерленд 1968. године када је креирао уређај познат као \textit{„Мач Дамокла“} (Sword of Damocles). Овај уређај је био толико тежак да је морао бити објешен о плафон, а графика се састојала од врло простих жичаних 3D модела генерисаних на рачунару.

Током осамдесетих година, термин „виртуелна реалност“ званично популаризује Џарон Ланир, оснивач компаније \textit{VPL Research}, која је развила прве комерцијалне рукавице (DataGlove) и наочаре за VR. Нажалост, тадашњи хардвер је био далеко испод потребних перформанси, што је довело до финансијског колапса и дугог периода познатог као „VR зима“. 

Тек са развојем паметних телефона, микросензора (акцелерометара и жироскопа) и изузетно моћних графичких чипова, долази до великог препорода 2012. године покретањем \textit{Kickstarter} кампање за \textcolor{blue}{Oculus Rift}. Од тада, технологија доживљава експлозиван раст, ширећи се изван оквира видео-игара у све сфере друштва.

\section{Теоријске основе и математичко моделовање}

Да би виртуелни свијет дјеловао увјерљиво и спријечио појаву дезоријентације код корисника, неопходно је ријешити низ сложених математичких задатака у реалном времену ($90+$ фрејмова у секунди).

\subsection{Математичке трансформације у 3D простору}

У процесу рендеровања, објекти се дефинишу у локалним координатним системима, а затим трансформишу у свјетске и камерине координате. Свака тачка $P(x, y, z)$ у простору се пресликава помоћу хомогених координата и матричног множења:
\begin{equation}
    P' = M_{\text{proj}} \cdot M_{\text{view}} \cdot M_{\text{model}} \cdot P
\end{equation}
гдје су $M_{\text{model}}$, $M_{\text{view}}$ и $M_{\text{proj}}$ матрице модела, погледа и пројекције. У физичким симулацијама често користимо и опште законе физике ради увођења реализма, попут познате Ајнштајнове једначине за еквиваленцију масе и енергије $E = mc^2$ код напредних симулација честица и експлозија у виртуелним свјетовима.

\subsection{Формалне дефиниције, леме и теореме}

\begin{definicija}
Виртуелна реалност (VR) је интерактивно компјутерски генерисано окружење које користи хардверске сензоре и дисплеје за стварање осјећаја потпуног уроњења (енгл. immersion) у вјештачки простор.
\end{definicija}

\begin{lema}
За постизање течног кретања без визуализационих трзаја, укупно кашњење система (енгл. motion-to-photon latency) мора бити стриктно мање од прага од $20 \text{ ms}$.
\end{lema}

\begin{teorema}
Ако је фреквенција освежавања екрана $f \ge 90 \text{ Hz}$ и ако је укупно кашњење система мање од критичне границе од $20 \text{ ms}$, онда људски мозак прихвата визуелну стимулацију као континуалну реалност без изазивања акутне кинетозе (мучнине кретања).
\end{teorema}

\section{Оптички системи и перцепција}

Људска перцепција простора почива на стереоскопији — чињеници да наше очи посматрају свијет са незнатно различитих тачака гледишта, што мозгу омогућава рачунање дубине.

\subsection{Оптички елементи и сочива}
Рани системи користили су обична сочива са значајним изобличењима. Савремени уређаји користе:
\begin{itemize}
    \item \textbf{Фреснелова сочива:} Омогућавају смањење тежине и дебљине оптичког склопа коришћењем концентричних прстенова.
    \item \textit{Pancake сочива:} Најновија технологија која пребија свјетлост више пута унутар оптичког пута, чинећи уређаје знатно тањим и ергономичнијим.
\end{itemize}

\subsection{Изазови у перцепцији}
Један од највећих проблема у оптичком дизајну јесте \textbf{конфликт између вергенце и акомодације}. У стварном свијету, наше очи се фокусирају на одређену удаљеност (акомодација) док се истовремено спајају на тачку интереса (вергенца). У VR уређајима, очи увијек фокусирају физички близак екран (око 2 метра или мање), док мозгу шаљемо сигнал да се објекат налази миљама далеко, што може довести до дуготрајног замора очију.

\section{Хардверске архитектуре и класификација уређаја}

Тржиште VR уређаја данас је диверзификовано према степену мобилности и рачунарске снаге коју користе.

\subsection{Класификација система}
\begin{enumerate}
    \item \textit{PC VR системи:} Зависе од екстерног рачунара високих перформанси (повезани каблом или бежично). Нуде најквалитетнију графику и најпрецизније праћење.
    \item \textit{Standalone (самостални) уређаји:} Садрже све компоненте (процесор, батерију, меморију) у самом кућишту наочара. Изузетно су практични због одсуства каблова.
    \item \textbf{Mobile VR системи:} Историјски значајна категорија гдје се паметни телефон користио као екран и процесор (нпр. Samsung Gear VR).
\end{enumerate}

\subsection{Упоредни преглед спецификација водећих модела}

\begin{table}[h]
\centering
\caption{Упоредна анализа техничких карактеристика актуелних VR уређаја}
\vspace{0.2cm}
\begin{tabular}{|l|c|c|c|}
\hline
\textbf{Модел уређаја} & \textbf{Резолуција по оку} & \textbf{Освежавање} & \textbf{Тип повезивања} \\ \hline
Meta Quest 3 & $2064 \times 2208$ пиксела & $120 \text{ Hz}$ & Самостални / PC \\ \hline
Valve Index & $1440 \times 1600$ пиксела & $144 \text{ Hz}$ & Жичани (DisplayPort) \\ \hline
Apple Vision Pro & $3664 \times 3140$ піксела & $100 \text{ Hz}$ & Самостални / Apple ecosystem \\ \hline
\end{tabular}
\end{table}

\section{Праћење положаја и интеракција}

Интеракција у виртуелном простору захтева прецизно праћење и покрета главе и положаја руку.

\subsection{Технологије праћења (Tracking)}
Постоје два основна приступа праћењу простора:
\begin{itemize}
    \item \textbf{Outside-in праћење:} Екстерне базне станице или камере постављене у соби прате свјетлеће или инфрацрвене маркере на наочарима и контролерима.
    \item \textit{Inside-out праћење:} Камере уграђене директно на уређај посматрају околину и користе алгоритме рачунарског вида (\textit{SLAM - Simultaneous Localization and Mapping}) ради креирања мапе простора у реалном времену.
\end{itemize}

\section{Примена виртуелне реалности у различитим дисциплинама}

Иако је првобитно везан за забаву, VR данас има кључну улогу у бројним индустријама.

\subsection{Медицина и хирургија}
У медицинском сектору, VR омогућава студентима медицине да виртуелно изводе сложене операције без ризика по стварне пацијенте. Такође, користи се у психотерапији за \textcolor{red}{третман различитих фобија и посттрауматичног стресног поремећаја} кроз контролисану експозицију пацијента изазовним ситуацијама.

\subsection{Инжењеринг и архитектура}
Инжењери користе VR за испитивање 3D прототипова аутомобила, авиона или зграда прије него што се покрене скупи физички процес производње или градње, чиме се штеде милиони евра.

\section{Изазови и будућност технологије}

Упркос невјероватном напретку, пред индустријом стоје озбиљни изазови. Главне препреке укључују \textbf{физичку тежину уређаја}, трајање батерије код самосталних уређаја, као и потпуно елиминисање мучнине кретања код осетљивијих корисника. Будућност доноси дубљу интеграцију са вјештачком интелигенцијом, развој лакших наочара са проширеном реалношћу (AR/MR), па чак и експерименте са интерфејсима мозак-рачунар (BCI).

\section{Закључак}

Технологија виртуелне реалности је прешла дуг пут од гломазних лабораторијских експеримената до софистицираних потрошачких уређаја који мијењају начин на који живимо, радимо и комуницирамо. \textit{Иако технички и перцептивни изазови и даље постоје}, несумњиво је да VR више није само научна фантастика, већ темељна технолошка платформа будућности која ће трајно редефинисати границе људског искуства.

\end{document}